% =====================================================================
% Capítulo 1 — Esboço de seção/interlúdio
% Sugestão de inserção: imediatamente após o parágrafo que termina com
% "[...] fazer alianças \parencite{stengers2010cosmopolitics1}." (tese
% como inscrição tecnocientífica), e antes do parágrafo que abre com
% "De certo modo, em sintonia com as críticas de Haraway e Stengers [...]".
%
% Observação técnica: para o build do LaTeX encontrar a figura, copie
%   rede_textual/infranodus_cap1_focus.png
% para
%   figuras/cap.1/infranodus_cap1_focus.png
% (mantendo a convenção do capítulo). O caminho abaixo já assume essa
% localização. Se preferir manter a figura em `rede_textual/`, basta ajustar
% o argumento de \includegraphics.
% =====================================================================

\subsection*{Interlúdio: a tese como inscrição que se mede a si mesma}

Se esta tese é uma inscrição tecnocientífica, então ela mesma pode ser
submetida ao tipo de tratamento que ela aplica aos seus interlocutores:
pode ser convertida em traço, em diagrama, em superfície combinável.
\textcite{Latour1986Visualization} sugere que a força de uma inscrição
está em ser \textit{móvel imutável}, isto é, transportável e
recombinável em uma única superfície plana onde inscrições heterogêneas
passam a ser comparadas. Os centros de cálculo, segue Latour, agem à
distância porque acumulam essas superfícies. O que ofereço aqui é um
gesto experimental e propositadamente recursivo: tomo o texto deste
capítulo como objeto e o submeto a uma \textbf{análise de rede textual}
(\textit{text network analysis}), método em que o texto é representado
como grafo de co-ocorrência de termos e as comunidades desse grafo são
lidas como tópicos latentes. A formulação contemporânea do método
deriva da tradição mais ampla das \textit{redes de co-ocorrência
lexical} \parencite{ferrercancho2001small} e da \textit{análise de
redes semânticas} \parencite{diesner2005revealing}, e foi sistematizada
por \textcite{paranyushkin2019infranodus} no \textit{pipeline} que aqui
reproduzo.\footnote{A análise foi conduzida com um \textit{script} em
Python que reproduz a metodologia descrita por
\textcite{paranyushkin2019infranodus}: o texto do capítulo, depurado
dos comandos \LaTeX{} e com as notas de rodapé reincorporadas ao
corpus, foi tokenizado em \textit{lemas} e percorrido por uma janela
deslizante de quatro \textit{tokens}; cada par de termos co-ocorrentes
na janela recebeu um peso decrescente com a distância (3, 2, 1). O
grafo ponderado foi podado para os 180 termos mais frequentes e para
arestas com peso $\geq 2$; em seguida, foram detectadas comunidades
pelo algoritmo de Louvain ponderado \parencite{blondel2008fast} e
calculadas medidas de \textit{degree} ponderado e \textit{betweenness
centrality}. Os artefatos da análise estão depositados em
\texttt{rede_textual/} no repositório desta tese (\texttt{.gexf} para
\textit{Gephi}, \texttt{.png}, \texttt{.csv}, e relatório
\texttt{.md}).}

Antes de ler os resultados, convém que eu explicite em que sentido as
figuras que se seguem e o grafo que acompanha esta tese no site são a
mesma análise em dois regimes. O procedimento tem duas camadas
sobrepostas. A primeira é a rede de co-ocorrência, que converte o texto
em nós, os termos, e em arestas, a proximidade entre eles, e da qual
extraio as medidas que leio adiante: o \textit{degree}, que mede a
centralidade de cada termo, o \textit{betweenness}, que localiza os
termos-ponte por onde um vocabulário transita para outro, e o NPMI, que
revela os pares mais colados. A segunda camada é a detecção de
comunidades pelo algoritmo de Louvain, que roda sobre esse mesmo grafo e
o reparte nos agrupamentos que a cor das figuras exibe, lidos aqui como
tópicos latentes. As duas camadas operam juntas nos dois lugares em que
mobilizo a análise, e a diferença entre eles está no escopo e na mídia,
com o método mantido idêntico. No corpo impresso da tese aplico o
procedimento capítulo a capítulo, rodando o Louvain dentro de cada um e
congelando o resultado em uma figura estática, uma inscrição no sentido
que \textcite{Latour1986Visualization} dá ao termo, um móvel imutável que
curo, legendo e faço entrar na leitura linear do argumento, e que me
devolve a estrutura temática interna daquele capítulo. No site que
acompanha a tese aplico o mesmo procedimento à tese inteira de uma vez,
rodando o Louvain sobre o conjunto dos capítulos concatenados, e o
apresento como um grafo navegável em que percorro os termos, abro as
passagens que os sustentam e alterno entre a tese como um todo e cada
capítulo. A escolha da mídia é ela própria parte do método não inocente
que assumo desde o capítulo~\ref{capitulo1}, pois a figura fixa faz uma
afirmação e a superfície navegável abre uma exploração, e cada uma
inscreve, à sua maneira, o gesto reflexivo de submeter esta tese ao mesmo
tratamento que ela dá aos seus objetos.\footnote{A diferença de escopo tem
uma consequência que registro para não confundi-la com ruído. Rodado sobre
a tese inteira, o Louvain produz uma única partição, em que cada termo
recebe uma comunidade válida para o todo e os agrupamentos são temas
transversais aos capítulos; rodado capítulo a capítulo, produz uma
partição independente para cada um, de sorte que um mesmo termo pode
pertencer a comunidades distintas em capítulos distintos, conforme a sua
vizinhança lexical se desloca. As comunidades por capítulo pedem, por
isso, que eu as rotule antes de compará-las, ao passo que a partição da
tese inteira já as oferece comparáveis. Preservo os dois regimes porque
cada um responde a uma pergunta: a estrutura interna de cada capítulo, de
um lado, e os temas que atravessam a tese inteira, de outro.}

A leitura desse grafo é, ela própria, uma operação de tradução: cada
nó é um termo, cada aresta é uma co-ocorrência, e cada cor é uma
comunidade detectada por um algoritmo cuja lógica é tão alheia à
hermenêutica antropológica quanto as redes neurais que descrevo nos
capítulos~\ref{capitulo3} e~\ref{capitulo4}. E, ainda assim, ele me
devolve algo. Os termos de maior \textit{degree} ponderado neste
capítulo são, em ordem, \textit{rede}, \textit{pesquisa},
\textit{etnografia}, \textit{artificial}, \textit{inteligência},
\textit{ciência}, \textit{latour}, \textit{campo}, \textit{método} e
\textit{objeto}. Os termos que mais costuram comunidades distintas
(maior \textit{betweenness}) são \textit{rede} (0,41), \textit{pesquisa}
(0,29) e \textit{etnografia} (0,19), seguidos de \textit{latour} e de
\textit{corte}: os três primeiros são os termos genéricos que permitem o
trânsito entre o vocabulário ator-rede, o vocabulário das ciências
sociais e o vocabulário da inteligência artificial, ao passo que
\textit{corte}, o corte agencial, comparece como a única ponte teórica
densa. A própria condição de \textit{ponte}, no sentido que
\textcite{Latour1986Visualization} dá ao termo, coincide aqui com a
condição de \textit{tradutor} no sentido latouriano mais amplo
\parencite{Latour2011CienciaEmAcao}.

A detecção de comunidades estabilizou os tópicos
latentes:\footnote{Os rótulos abaixo correspondem aos termos de maior
\textit{degree} dentro de cada comunidade; preservo a forma lematizada
produzida pelo \textit{script} para evitar tradução interpretativa
precoce.}
(i)~o \textbf{método e a teoria ator-rede}, a maior comunidade, em torno
de \textit{latour}, \textit{método}, \textit{corte}, \textit{strathern}
e \textit{gesto}, com um segundo agrupamento ator-rede em \textit{rede},
\textit{análise}, \textit{ator} e \textit{teoria};
(ii)~a \textbf{etnografia e o campo} em torno de \textit{etnografia},
\textit{descrição}, \textit{prática} e \textit{materiais};
(iii)~a \textbf{inteligência artificial como objeto} em torno de
\textit{pesquisa}, \textit{artificial}, \textit{inteligência},
\textit{objeto} e \textit{modelo};
(iv)~as \textbf{ciências sociais} em torno de \textit{ciência},
\textit{campo}, \textit{dado}, \textit{sociais} e \textit{escrita};
(v)~a \textbf{simetria humano-máquina} em torno de \textit{humano},
\textit{relação}, \textit{máquina}, \textit{parcial} e \textit{plano};
(vi)~um polo teórico denso de Law e Mol em torno de \textit{hinterland},
\textit{otherness}, \textit{ausência}, \textit{presença} e
\textit{manifesta}.

\begin{figure}[!htb]
\centering
\includegraphics[width=1.0\textwidth,keepaspectratio]{figuras/cap.1/infranodus_cap1_focus.png}
\caption[Rede textual do capítulo~\ref{capitulo1} (\textit{text network
analysis})]{Rede textual do capítulo~\ref{capitulo1}
(\textit{text network analysis}). Os nós são os 100 termos de maior
\textit{degree} ponderado; o tamanho do nó acompanha o
\textit{degree}; a cor indica a comunidade detectada por Louvain
ponderado; as arestas, arqueadas, recebem a mistura das cores das
comunidades que ligam (peso~$\geq 3$). Versão de alta resolução,
arquivo \texttt{.gexf} para o \textit{Gephi} e relatório completo
em \texttt{rede_textual/} no repositório desta tese.}
\label{fig:infranodus-cap1}
\end{figure}

O que essa cartografia me dá tem a forma de um conjunto de
\textit{chamadas de atenção}. A primeira é uma confirmação modesta:
os termos centrais do grafo coincidem com os conceitos que mobilizo
explicitamente, e os pontos de maior intermediação são justamente os
termos mais genéricos do capítulo (\textit{rede}, \textit{pesquisa},
\textit{etnografia}), o que sugere que a costura argumentativa vem sendo
feita por palavras-coringa que ainda esperam ser teorizadas com mais
força. A segunda é uma tensão que o grafo expõe: as ligações ponderadas
mais fracas caem entre a comunidade do \textbf{método e da teoria
ator-rede} (\textit{método}, \textit{latour}, \textit{corte}) e a da
\textbf{inteligência artificial como objeto} (\textit{artificial},
\textit{inteligência}, \textit{objeto}), e também entre a
\textbf{etnografia} e a \textbf{simetria humano-máquina}
(\textit{humano}, \textit{parcial}, \textit{agência}). O vocabulário com
que prometo descrever e o objeto técnico que me proponho a descrever
ainda correm em paralelo. A formulação de \textcite{Law2004After} sobre
o \textit{method assemblage} ajuda a nomear o que isso quer dizer: a
\textit{otherness} apagada para que a presença se sustente é, neste
capítulo, a articulação entre esses registros, que ainda preciso fazer
trabalhar juntos. A terceira chamada diz respeito ao polo teórico denso
de Law e Mol (\textit{otherness}, \textit{ausência}, \textit{presença},
\textit{manifesta}), que os pares NPMI mais fortes revelam bem costurado
por dentro, \textit{ausência}~$\leftrightarrow$~\textit{manifesta} (0,87)
e \textit{presença}~$\leftrightarrow$~\textit{ausência} (0,70), e que no
entanto pende da rede maior por poucos fios: uma promessa teórica que o
capítulo abre e que os capítulos subsequentes precisam costurar ao todo,
sob risco de restar como nota de rodapé errante.

Reconheço o paradoxo do gesto. Aplicar análise de rede textual a um
capítulo que invoca a teoria ator-rede é arriscar tomar a metáfora
pelo método; e a similaridade entre os dois usos da palavra
\textit{rede} é superficial: a rede ator-rede é uma cadeia de
associações heterogêneas e ontologicamente plana, ao passo que a rede
textual aqui gerada é um grafo bidimensional de co-ocorrências de
\textit{tokens} ponderado por proximidade.\footnote{A confusão entre
os dois sentidos de \textit{rede} é discutida por
\textcite[pp.~129--132]{latour2005reassembling}, que distingue
\textit{rede} como ferramenta de inscrição e como diagnóstico
ontológico. Mantenho a distinção: o que a análise de rede textual
produz é uma inscrição; o que esta tese descreve é uma rede
sociotécnica. As duas se encontram apenas no gesto de inscrever uma
na outra.} Talvez seja exatamente isso que justifica
o gesto: a rede textual se acrescenta à etnografia como mais uma
inscrição que esta tese mobiliza, no mesmo registro das imagens, dos
diagramas e do site que a acompanham. Como qualquer inscrição, ela
ordena, hierarquiza, inclui e exclui, ou seja, faz política ontológica
\parencite[p.~143]{Law2004After}. E talvez essa seja a versão mais
honesta do gesto reflexivo que esta tese tenta sustentar: medir o
próprio texto pelos mesmos instrumentos com que se compromete a medir o
C4AI, sabendo que medir é também produzir aquilo que se mede.

% =====================================================================
% Sugestões de entradas bibliográficas (caso ainda não estejam no .bib):
%
% @inproceedings{paranyushkin2019infranodus,
%   author    = {Paranyushkin, Dmitry},
%   title     = {{InfraNodus}: Generating Insight Using Text Network Analysis},
%   booktitle = {The World Wide Web Conference (WWW '19)},
%   year      = {2019},
%   pages     = {3584--3589},
%   publisher = {ACM},
%   doi       = {10.1145/3308558.3314123}
% }
%
% @article{ferrercancho2001small,
%   author  = {Ferrer i Cancho, Ramon and Sol{\'e}, Ricard V.},
%   title   = {The small world of human language},
%   journal = {Proceedings of the Royal Society B: Biological Sciences},
%   volume  = {268},
%   number  = {1482},
%   pages   = {2261--2265},
%   year    = {2001},
%   doi     = {10.1098/rspb.2001.1800}
% }
%
% @incollection{diesner2005revealing,
%   author    = {Diesner, Jana and Carley, Kathleen M.},
%   title     = {Revealing Social Structure from Texts: Meta-Matrix
%                Text Analysis as a Novel Method for Network Text Analysis},
%   booktitle = {Causal Mapping for Information Systems and Technology Research:
%                Approaches, Advances, and Illustrations},
%   editor    = {Narayanan, V. K. and Armstrong, D. J.},
%   publisher = {Idea Group Publishing},
%   year      = {2005},
%   pages     = {81--108}
% }
%
% @article{blondel2008fast,
%   author  = {Blondel, Vincent D. and Guillaume, Jean-Loup and
%              Lambiotte, Renaud and Lefebvre, Etienne},
%   title   = {Fast unfolding of communities in large networks},
%   journal = {Journal of Statistical Mechanics: Theory and Experiment},
%   volume  = {2008},
%   number  = {10},
%   pages   = {P10008},
%   year    = {2008},
%   doi     = {10.1088/1742-5468/2008/10/P10008}
% }
% =====================================================================
